\begin{lemma}\label{lem:dyadic-cuspidal-diamond}
Let $N\geq5$ be odd, put $X=X_1(N)$ and $J=J_1(N)$, and let
$d\geq1$. Suppose that the following conditions hold:
\begin{enumerate}
\item $Y_1(N)(\mathbb F_{2^f})=\emptyset$ for $1\leq f\leq d$.
\item Every closed cusp of $X_{\mathbb F_2}$ of degree at most $d$
is the reduction of a rational cuspidal divisor of the same degree.
\item There is a diamond $h=\langle a\rangle$ of order $r$ such that
$(h-1)J(\mathbb Q)$ is finite and $r\nmid d$.
\item $\operatorname{gon}_{\mathbb C}(X)>4d$.
\end{enumerate}
Then there is no effective rational divisor of degree $d$ on $X$
with entirely noncuspidal support. In particular, no elliptic curve
over a degree-$d$ field has a point of order $N$.
\end{lemma}
\begin{proof}
Suppose that $D$ is such a divisor.
By (1), its reduction at $2$ is entirely cuspidal.
By (2), it has a rational cuspidal lift $C$ of degree $d$.
The point $y=(h-1)[D-C]$ is rational torsion and reduces to zero.
The torsion kernel of reduction at $2$ is killed by $2$; see the
proof of \cite[Proposition~3.4]{DKSS}.
Thus
\[
 2hD+2C\sim2D+2hC.
\]
Both sides have degree $4d$.
By (4), they are equal as divisors.
Separating their cuspidal and noncuspidal parts gives $hD=D$.

Every point in the support of $D$ has cuspidal reduction, so its
$j$-invariant has negative valuation at a place above $2$.
A CM $j$-invariant is an algebraic integer. Hence these points are
non-CM, and every geometric orbit under $h$ has length $r$.
It follows that $r\mid\deg D=d$, contradicting (3).
\end{proof}
